\subsection{Bài tập tự luận}
%%==========Bài 1
\begin{bt}
	Sử dụng MTCT tính các tỉ số lượng giác và làm tròn kết quả đến chữ số thập phân thứ ba:
	\begin{listEX}[2]
	\item
	$\sin 27^\circ$, 
	$\cos 32^\circ 15'$, 
	$\tan 52^\circ 12'$ 
	và $\cot 35^\circ 23'$.
	\item
	$\sin 40^\circ 54'$,
	$\cos 52^\circ 15'$,
	$\tan 69^\circ 36'$
	và $\cot 25^\circ 18'$.
	\end{listEX}
	\loigiai{
%	\begin{listEX}
%	\item
%	\begin{tabular}{|l|l|l|}
%	\hline
%	Để tính & Bấm phím & Kết quả \\ \hline
%	$\sin 27^\circ$ & \sink \key{2} \key{7} \key{x} \key{=} & $0{,}4539904997$ \\ \hline
%	$\cos 32^\circ 15'$ & \cosk \key{3} \key{2} \key{x} \key{1} \key{5} \key{x} \key{=} & $0{,}8457278217$ \\ \hline
%	$\tan 52^\circ 12'$ & \tank \key{5} \key{2} \key{x} \key{1} \key{2} \key{x} \key{=} & $1{,}289192232$ \\ \hline
%	$\cot 35^\circ 23'$ & \tank \key{3} \key{5} \key{x} \key{2} \key{3} \key{x} \key{=} \key{u} \key{=} & $1{,}408003909$ \\ \hline
%	\end{tabular} \\
%	Làm tròn đến chữ số thập phân thứ ba ta được $\sin 27^\circ \approx 0{,}454$; $\cos 32^\circ 15'\approx 0{,}846$; $\tan 52^\circ 12' \approx 1{,}289$; $\cot 35^\circ 23' \approx 1{,}408$.\\
%	Lưu ý, $\cot 35^\circ 23' = \dfrac{1}{\tan 35^\circ 23'}$.
%	\item
%	\begin{tabular}{|l|l|l|}
%	\hline
%	Để tính & Bấm phím & Kết quả \\ \hline
%	$\sin 40^\circ 54'$ & \sink \key{4} \key{0} \key{x} \key{5} \key{4} \key{x} \key{=} & $0{,}6547408137$ \\ \hline
%	$\cos 52^\circ 15'$ & \cosk \key{5} \key{2} \key{x} \key{1} \key{5} \key{x} \key{=} & $0{,}61221728$ \\ \hline
%	$\tan 69^\circ 36'$ & \tank \key{6} \key{9} \key{x} \key{3} \key{6} \key{x} \key{=} & $2{,}688918967$ \\ \hline
%	$\cot 25^\circ 18'$ & \tank \key{2} \key{5} \key{x} \key{1} \key{8} \key{x} \key{=} \key{u} \key{=} & $2{,}115516356$ \\ \hline
%	\end{tabular} \\
%	Làm tròn đến chữ số thập phân thứ ba ta được $\sin 40^\circ 54' \approx 0{,}655$; $\cos \cos 52^\circ 15' \approx 0{,}612$; $\tan 69^\circ 36' \approx 2{,}689$; $\cot 25^\circ 18' \approx 2{,}116$.
%	\end{listEX}
	}
\end{bt}
%%==========Bài 2
\begin{bt}
	Dùng MTCT, tìm các góc (làm tròn đến phút) biết
	\begin{listEX}
	\item 
	$\sin \alpha _1 = 0{,}3214$, 
	$\cos \alpha _2 =0{,}4321$, 
	$\tan \alpha _3 = 1{,}2742$ 
	và $\cot \alpha _4 = 1{,}5384$.
	\item
	$\sin \alpha_1 = 0{,}3782$,
	$\cos \alpha_1 = 0{,}6251$,
	$\tan \alpha_1 = 2{,}154$
	và $\cot \alpha_1 = 3{,}253$.
	\end{listEX}
	\loigiai{
%	\begin{listEX}
%	\item
%	\begin{tabular}{|l|l|l|l|}
%	\hline
%	Biết & Bấm phím & Kết quả & Bấm tiếp \key{x} \\ \hline
%	$\sin \alpha _1 = 0{,}3214$ & \key{q} \sink \key{0} \key{.} \key{3} \key{2} \key{1} \key{4} \key{=} & $18{,}74761209$ & $18^\circ 44' 51{,}4"$ \\ \hline
%	$\cos \alpha _2 =0{,}4321$ & \key{q} \cosk \key{0} \key{.} \key{4} \key{3} \key{2} \key{1} \key{=} & $64{,}39909458$ &$64^\circ 23' 56{,}74"$ \\ \hline
%	$\tan \alpha _3 = 1{,}2742$ & \key{q} \tank \key{1} \key{.} \key{2} \key{7} \key{4} \key{2} \key{=} & $51{,}87495892$ & $51^\circ 52' 29{,}85"$ \\ \hline
%	$\cot \alpha _4 = 1{,}5384$ & \key{q} \tank \key{1} \key{.} \key{6} \key{3} \key{8} \key{4} \key{u} \key{=} & $33{,}02491482$ & $33^\circ 1' 29{,}69"$ \\ \hline
%	\end{tabular} \\
%	Làm tròn đến phút ta được: $\alpha _1 \approx 18^\circ 45'$, $\alpha _2 \approx 64^\circ 24'$, $\alpha _3 \approx 51^\circ 52'$, $\alpha _4 \approx 33^\circ 1'$.
%	\item
%	\begin{tabular}{|l|l|l|l|}
%	\hline
%	Biết & Bấm phím & Kết quả & Bấm tiếp \key{x} \\ \hline
%	$\sin \alpha = 0{,}3782$ & \key{q} \sink \key{0} \key{.} \key{3} \key{7} \key{8} \key{2} \key{=} & $22{,}222231$ & $22^\circ 13' 20{,}03"$ \\ \hline
%	$\cos \alpha = 0{,}6251$ & \key{q} \cosk \key{0} \key{.} \key{6} \key{2} \key{5} \key{1} \key{=} & $51{,}31047244$ &$51^\circ 18' 37{,}7"$ \\ \hline
%	$\tan \alpha = 2{,}154$ & \key{q} \tank \key{2} \key{.} \key{1} \key{5} \key{4} \key{=} & $65{,}09679426$ & $65^\circ 5' 48{,}46"$ \\ \hline
%	$\cot \alpha = 3{,}253$ & \key{q} \tank \key{3} \key{.} \key{2} \key{5} \key{3} \key{u} \key{=} & $17{,}08787556$ & $17^\circ 5' 16{,}35"$ \\ \hline
%	\end{tabular} \\
%	Làm tròn đến phút ta được:
%	$\alpha_1 \approx 22^\circ 13'$,
%	$\alpha_2 \approx 51^\circ 19'$,
%	$\alpha_3 \approx 65^\circ 6'$,
%	$\alpha_4 \approx 17^\circ 5'$.
%	\end{listEX}
	}
\end{bt}
%%==========Bài 3
\begin{bt}
	Tính các tỉ số lượng giác của góc $\widehat{ABC}=\alpha$ trong tam giác $ABC$ vuông tại $A$ có $AB=12, AC=9, BC=15$
	\loigiai{
	Xét tam giác $ABC$ vuông tại $A$, $\widehat{B}=\alpha$, ta có\\
	$\sin\alpha=\dfrac{AC}{BC}=\dfrac{9}{15}=0{,}6$;\qquad $\cos\alpha=\dfrac{AB}{BC}=\dfrac{12}{15}=0{,}8$;\\
	$\tan\alpha=\dfrac{AC}{AB}=\dfrac{9}{12}=0{,}75$;\qquad $\cot\alpha=\dfrac{AB}{AC}=\dfrac{12}{9}=\dfrac{4}{3}$.
	}
\end{bt}
%%==========Bài 4
\begin{bt}
	Tính các tỉ số lượng giác của góc nhọn $A$ trong mỗi tam giác vuông $ABC$ có $\widehat{B}=90^\circ$ ở hình sau.
	\begin{center}
	\begin{tikzpicture}[declare function={a=2.5;},line join=round,line cap=round,font=\footnotesize,scale=1]
	\path (0:0) coordinate (B)
	(0:a) coordinate (C)
	(B)+(90:0.7*a) coordinate (A)
	($(B)!0.5!(C)$) coordinate (D) node[below,shift={(-90:0.6cm)}] {a)}
	;
	\path pic[draw,angle radius=5pt]{right angle= A--B--C};
	\path pic[draw,angle radius=7pt]{angle= B--A--C};
	\draw (A)--(B) node[midway,left] {$3$}
	--(C) node[midway,below] {$4$}
	--cycle node[midway,right,yshift=3pt] {$5$}
	;
	\foreach \t/\g in {A/90,B/-90,C/-90}{
	\path (\t) node[shift={(\g:7pt)},font=\scriptsize]{$ \t $};
	}
	\end{tikzpicture} \hspace*{0.6cm}
	\begin{tikzpicture}[declare function={a=2;},line join=round,line cap=round,font=\footnotesize,scale=1]
	\path 
	(0:a) coordinate (C)
	(180:a) coordinate (B)
	(150:a) coordinate (A)
	($(B)!0.5!(C)$) coordinate (D) node[shift={(-90:1cm)}] {b)}
	;
	\path pic[angle eccentricity=2,draw,angle radius=12pt]{angle= A--C--B};
	\path pic[draw,angle radius=5pt]{right angle= B--A--C};
	\draw 
	(A)--(B) node[midway,left] {$1$}
	--(C) node[midway, below] {$\sqrt{17}$}--cycle node[midway, above] {$4$}
	;
	\foreach \t/\g/\i in {A/90/B,B/-90/C,C/-90/A}{
	\path (\t) node[shift={(\g:7pt)}]{$ \i $};
	}
	\end{tikzpicture}
	\begin{tikzpicture}[declare function={a=2.5;},line join=round,line cap=round,font=\footnotesize,scale=1]
	\path (0:0) coordinate (B)
	(0:a) coordinate (C)
	(C)+(-90:0.85*a) coordinate (A)
	(A)+(180:0.5*a) coordinate (O) node[shift={(-90:10pt)}] {c)}
	;
	\path pic[draw,angle radius=5pt]{right angle= A--C--B};
	\path pic[draw,angle radius=7pt]{angle= C--A--B};
	\draw (A)--(B) node[midway,left,xshift=-3pt] {$3$}
	--(C) 
	--cycle node[midway,right,yshift=4pt] {$2$}
	;
	\foreach \t/\g/\i in {A/-90/A,B/90/C,C/90/B}{
	\path (\t) node[shift={(\g:7pt)}]{$ \i $};
	}
	\end{tikzpicture} \hspace*{0.6cm}
	\begin{tikzpicture}[declare function={a=2.5;},line join=round,line cap=round,font=\footnotesize,scale=1]
	\path (0:0) coordinate (B)
	(180:a) coordinate (C)
	(B)+(90:0.7*a) coordinate (A)
	($(B)!0.5!(C)$) coordinate (D) node[below,shift={(-90:0.6cm)}] {d)}
	;
	\path pic[draw,angle radius=5pt]{right angle= A--B--C};
	\path pic[draw,angle radius=7pt]{angle= B--C--A};
	\draw (A)--(B) node[midway,right] {$\sqrt{6}$}
	--(C) node[midway,below] {$\sqrt{10}$}
	--cycle 
	;
	\foreach \t/\g/\i in {A/90/C,B/-90/B,C/-90/A}{
	\path (\t) node[shift={(\g:7pt)}]{$ \i $};
	}
	\end{tikzpicture}
	\end{center}
	\loigiai{
	Xét $\triangle ABC$ vuông tại $B$ (Hình a), ta có
	\begin{enumEX}[\itemCI]{4}
	\item $\sin A=\dfrac{BC}{AC}=\dfrac{4}{5}$;
	\item $\cos A=\dfrac{AB}{AC}=\dfrac{3}{5}$;
	\item $\tan A=\dfrac{BC}{AB}=\dfrac{4}{3}$;
	\item $\cot A=\dfrac{AB}{BC}=\dfrac{3}{4}$.
	\end{enumEX}
	Xét $\triangle ABC$ vuông tại $B$ (Hình b), ta có
	\begin{enumEX}[\itemCI]{4}
	\item $\sin A=\dfrac{BC}{AC}=\dfrac{1}{\sqrt{17}}$;
	\item $\cos A=\dfrac{AB}{AC}=\dfrac{4}{\sqrt{17}}$;
	\item $\tan A=\dfrac{BC}{AB}=\dfrac{1}{4}$;
	\item $\cot A=\dfrac{AB}{BC}=\dfrac{4}{1}$.
	\end{enumEX}
	Xét $\triangle ABC$ vuông tại $B$ (Hình c), ta có	$BC=\sqrt{3^2-2^2}=\sqrt{5}$;
	\begin{enumEX}[\itemCI]{4}
	\item $\sin A=\dfrac{BC}{AC}=\dfrac{\sqrt{5}}{3}$;
	\item $\cos A=\dfrac{AB}{AC}=\dfrac{2}{3}$;
	\item $\tan A=\dfrac{BC}{AB}=\dfrac{\sqrt{5}}{2}$;
	\item $\cot A=\dfrac{AB}{BC}=\dfrac{2}{\sqrt{5}}$.
	\end{enumEX}
	Xét $\triangle ABC$ vuông tại $B$ (Hình d), ta có $AC=\sqrt{\left(\sqrt{10}\right)^2-\left(\sqrt{6}\right)^2}=2$;
	\begin{enumEX}[\itemCI]{4}
	\item $\sin A=\dfrac{BC}{AC}=\dfrac{\sqrt{6}}{2}$;
	\item $\cos A=\dfrac{AB}{AC}=\dfrac{\sqrt{10}}{2}$;
	\item $\tan A=\dfrac{BC}{AB}=\dfrac{\sqrt{6}}{\sqrt{10}}$;
	\item $\cot A=\dfrac{AB}{BC}=\dfrac{\sqrt{10}}{\sqrt{6}}$.
	\end{enumEX}
	}
\end{bt}
%==========Bài 5
\begin{bt} %Ví dụ 1
	Cho tam giác $ABC$ vuông tại $A$, có $AB=3$ cm, $AC=4$ cm. Hãy tính các tỉ số lượng giác $\sin \alpha$, $\cos \alpha$, $\tan \alpha$ với $\alpha = \widehat{B}$.
	\loigiai{
	\immini{
	Xét $\triangle ABC$ vuông tại $A$, $\widehat{B}=\alpha$.\\
	Theo định lí Pythagore, ta có: $BC^2=AC^2 + AB^2 = 4^2 +3^2 = 25$ nên $BC=5$ (cm).\\
	Theo định nghĩa của tỉ số lượng giác sin, côsin , tang, ta có\\
	$\sin \alpha = \dfrac{AC}{BC} = \dfrac{4}{5}$, $\cos \alpha = \dfrac{AB}{BC} = \dfrac{3}{5}$, $\tan \alpha = \dfrac{AC}{AB} = \dfrac{4}{3}$.
	}
	{
	\begin{tikzpicture}[>=stealth,line join=round,line cap=round,font=\footnotesize,scale=0.65]
		\path 
		(0,0) coordinate (B)
		(3,0) coordinate (A)
		(3,4) coordinate (C);
		\draw (0,0) node[left]{$B$}--(1.5,0) node[below]{$3$ cm}--(3,0) node[right]{$A$}--(3,2) node[right]{$4$ cm}--(3,4) node[above]{$C$}--(0,0);
		\pic[draw, angle radius=3mm]{right angle=B--A--C};
		\pic[draw, angle radius=5mm]{angle=A--B--C} node at ($(B)+(25:0.8)$){$\alpha$};
	\end{tikzpicture}
	}
	} 
\end{bt}
%%==========Bài 6
\begin{bt}
	\immini{Cho hình thoi $ABCD$ có hai đường chéo cắt nhau tại điểm $O$.
	\begin{listEX}
	\item Tỉ số $\dfrac{OB}{AB}$ là sin của góc nhọn nào? Tỉ số $\dfrac{OB}{BC}$ là côsin của góc nhọn nào?
	\item Viết tỉ số lượng giác của mỗi góc nhọn sau: $\tan \widehat{OCD}$, $\cot \widehat{OAD}$.
	\end{listEX}}
	{\begin{tikzpicture}[line join = round, line cap = round,>=stealth,font=\footnotesize,scale=0.7]
	\def \a{3};
	\path 
	(0,0) coordinate (O)
	(-\a,0)	coordinate (A)	
	(\a,0)coordinate (C)	
	($(O)!0.5!90:(C)$)coordinate (B)
	($(O)!0.5!-90:(C)$)coordinate (D)
	;
	\draw (A)--(B)--(C)--(D)--cycle (A)--(C) (B)--(D); 
	\foreach \x/\g in {O/45,A/180,B/90,C/0,D/-90}\fill[black] (\x) circle (1pt) ($(\x)+(\g:4mm)$) node{$\x$};
	\foreach \a/\b/\c in {C/O/A}{\draw pic[draw,angle radius=2mm] {right angle = \a--\b--\c};}	
	\end{tikzpicture}}
	\loigiai{
	Hình thoi $ABCD$ có hai đường chéo cắt nhau tại điểm $O$ nên $AC$ vuông góc với $BD$ tại $O$.
	\begin{listEX}[2]
	\item Tam giác $OAB$ vuông tại $O$ nên $\dfrac{OB}{AB}=\sin \widehat{OAB}$.\\
	Tam giác $OBC$ vuông tại $O$ nên $\dfrac{OB}{BC}=\cos \widehat{OBC}$.
	\item Tam giác $OCD$ vuông tại $O$ nên tan $\widehat{OCD}=\dfrac{OD}{OC}$.\\
	Tam giác $OA D$ vuông tại $O$ nên cot $\widehat{OAD}=\dfrac{OA}{OD}$.
	\end{listEX}
	}
\end{bt}
%%==========Bài 7
\begin{bt} %Luyện tập 1
	Cho tam giác $ABC$ vuông tại $A$ có $AB=5$ cm, $AC=12$ cm. Hãy tính các tỉ số lượng giác của góc $B$.
	\loigiai{
	\immini{
	Xét $\triangle ABC$ vuông tại $A$\\
	Theo định lí Pythagore, ta có \\
	$BC^2=AC^2 + AB^2 = 12^2 +5^2 = 169$ nên $BC=13$ (cm).\\
	Theo định nghĩa của tỉ số lượng giác sin, côsin , tang, côtang ta có\\
	$\sin \alpha = \dfrac{AC}{BC} = \dfrac{12}{13}$, $\cos \alpha = \dfrac{AB}{BC} = \dfrac{5}{13}$,\\
	$\tan \alpha = \dfrac{AC}{AB} = \dfrac{12}{5}$, $\cot \alpha = \dfrac{AB}{AC} = \dfrac{5}{12}$.
	}{
	\begin{tikzpicture}[>=stealth,line join=round,line cap=round,font=\footnotesize,xscale=0.5,yscale=0.25]
	\path 
	(0,0) coordinate (A)
	(5,0) coordinate (B)
	(0,12) coordinate (C);
	\draw (0,0) node[left]{$A$}--(2.5,0) node[below]{$5$ cm}--(5,0) node[right]{$B$}--(0,12) node[above]{$C$}--(0,6) node[left]{$12$ cm}--(0,0);
	\pic[draw, angle radius=3mm]{right angle=B--A--C};
	\pic[draw, angle radius=3mm]{angle=C--B--A};
	\end{tikzpicture}}}
\end{bt}
%%==========Bài 8
\begin{bt}%[9H1B2]
	$\triangle ABC$ vuông tại $A$ có $BC = 2AB$. Tính các tỉ số lượng giác của góc $C$.
	\loigiai
	{
	Ta đặt $AB = m$ thì $BC = 2m$, suy ra $AC^2 = BC^2 - AB^2 = 4m^2 - m^2 = 3m^2 \Rightarrow AC = m \sqrt{3}$.
	\immini
	{
	Ta có $\begin{aligned}[t]
	&\sin C =\dfrac{AB}{BC}=\dfrac{m}{2m}=\dfrac{1}{2}; \\
	& \cos C =\dfrac{AC}{BC}=\dfrac{m\sqrt{3}}{2m}=\dfrac{\sqrt{3}}{2}; \\
	& \tan C =\dfrac{AB}{AC}= \dfrac{m}{m \sqrt{3}}=\dfrac{1}{\sqrt{3}}; \\
	& \cot C = \dfrac{AC}{AB} = \dfrac{m \sqrt{3}}{m} = \sqrt{3}.
	\end{aligned}$
	}
	{
	\begin{tikzpicture}[>=stealth,line join=round,line cap=round,font=\footnotesize,scale=0.8]
	\tkzDefPoints{0/0/A}
	\coordinate (B) at ($(A)+(0,3)$);
	\coordinate (C) at ($(A)+(4,0)$);
	\pgfresetboundingbox
	\tkzMarkRightAngles[size=0.25](B,A,C)
	\tkzDrawSegments(A,B B,C C,A)
	\tkzDrawPoints[fill=black](A,B,C)
	%
	\pgfresetboundingbox
	\tkzMarkAngles[size=.5,arc=l](B,C,A)
	\tkzLabelPoints[above](B)
	\tkzLabelPoints[below](A,C)
	\end{tikzpicture}
	}
	}
\end{bt}
%%==========Bài 9
\begin{bt}%[9H1B2]
	Tam giác $ABC$ cân tại $A$, có $BC = 6$, đường cao $AH = 4$. Tính các tỉ số lượng giác của góc $B$. 
	\loigiai
	{
	Ta có $BH = 6 : 2 = 3$; $AB = \sqrt{4^2 + 3^2} = 5$.
	\immini
	{
	Do đó $\begin{aligned}[t]
	&\sin B =\dfrac{AH}{AB}=\dfrac{4}{5}= 0,8; \\
	& \cos B =\dfrac{BH}{AB}=\dfrac{3}{5}= 0,6;\\
	& \tan B = \dfrac{AH}{AB} = \dfrac{4}{3}; \\
	& \cot B = \dfrac{BH}{AH} = \dfrac{3}{4} = 0,75.
	\end{aligned}$
	}
	{
	\begin{tikzpicture}[>=stealth,line join=round,line cap=round,font=\footnotesize,scale=0.6]
	\tkzDefPoints{0/0/B}
	\coordinate (C) at ($(B)+(6,0)$);
	\tkzDefTriangle[two angles = 55 and 55](B,C) \tkzGetPoint{A}
	\tkzDefPointBy[projection = onto B--C](A) \tkzGetPoint{H}
	%
	\pgfresetboundingbox
	\tkzMarkRightAngles[size=0.25](A,H,C)
	\tkzDrawSegments(A,B B,C C,A A,H)
	\tkzDrawPoints[fill=black](A,B,C,H)
	\tkzLabelPoints[above](A)
	\tkzLabelPoints[above left](H)
	\tkzLabelPoints[left](B)
	\tkzLabelPoints[right](C)
	\draw[|<->|] ([yshift=-0.6]B)--([yshift=-0.6]C) node[below,midway,sloped]{$6$};
	%
	\path (A)--(H) node[right,midway]{$4$};
	\end{tikzpicture}
	}
	}
\end{bt}
%%==========Bài 10
\begin{bt}%[9H1B2]
Cho $\triangle ABC$ ($\widehat{BAC}=90^\circ$), $AH\perp BC$, biết $AB=6$, $BH=3$.	Tính $\tan C$ 
	\loigiai
	{
	Ta có $AH^2 = AB^2 - BH^2 = 6^2 - 3^2 = 27 \Rightarrow AH = 3 \sqrt{3}$.\\
	Do đó $\tan C = \cot B = \dfrac{BH}{AH} = \dfrac{3}{3 \sqrt{3}} = \dfrac{1}{\sqrt{3}}$.
	}
\end{bt}
%%==========Bài 11
\begin{bt}%[9H1B2]
	Cho tam giác $\triangle OMN$ vuông tại $O$ có $OH \perp MN$ tại $H$, $HM=1$, $HN=3$. Tính $\sin M + \cos N $ 
	\loigiai
	{
	Ta có $OH^2 = HM \cdot HN = 1 \cdot 3 = 3 \Rightarrow OH = \sqrt{3}$; $OM = \sqrt{1 + 3} = 2$.\\
	Do đó $\sin M = \dfrac{OH}{OM} = \dfrac{\sqrt{3}}{2}$.\\
	Mặt khác $\cos N = \sin M = \dfrac{\sqrt{3}}{2}$ nên $\sin M + \cos N = \sqrt{3}$.
	}
\end{bt}
%%==========Bài 12
\begin{bt}
	\immini{Tia nắng chiếu qua điểm $B$ của nó tòa nhà tạo với mặt đất một góc $x$ và tạo với cạnh $AB$ của tòa nhà một góc $y$ (Hình $9$). Cho biết $\cos x\approx 0{,}78$ và $\cot x\approx 1{,}25$. Tính $\sin y$ và $\tan y$ (kết quả làm tròn đến hàng phần trăm).}{
	\begin{tikzpicture}[>=stealth,font=\footnotesize, yscale=.3,xscale=0.4]
	\coordinate (O) at (0,0) ;
	\def\nhathap{5} %so tang nha 
	\def\kc{10} %khoang cach 2 nha
	\def\dai{1} %chieu dai 1 khoi
	\def\cao{2} %chieu cao 1 khoi
	\coordinate (H) at (\kc,0);
	\newcommand{\xaynha}[2]{
	\draw[very thick] (#1,#2)rectangle(#1+\dai,#2+\cao);
	\draw[very thick,fill=black] (#1,#2)rectangle(#1+\dai,#2+0.25*\cao);
	\draw[very thick] (#1+0.2*\dai,#2+0.25*\cao)rectangle(#1+0.8*\dai,#2+0.9*\cao);
	}
	% Xay nha
	\foreach \j in {1,...,\nhathap}{
	\foreach \i in {0,...,1}{
	\xaynha{\i*1.1}{\j*\cao}
	}
	}
	\draw[line width=0.1cm] (-0.1,\cao*\nhathap+\cao)--(2*\dai*1.1,\cao*\nhathap+\cao);
	\draw ($(0,2)+(2.1,0)$) coordinate (A)--($(90:12)+(2.1,0)$) coordinate (B)--(5,2)coordinate (C)--cycle
	;
	\foreach \x/\g in {A/-90,B/90,C/-90}\fill[red] (\x) circle (1pt)+(\g:6mm) node[black]{$ \x $};
	\path
	(A)--(C)%node[below,midway]{Mặt sàn}
	(A)--(C)
	(A)--(B)
	(B)++(30:1.25cm)%node{Sân thượng}
	;
	\path pic[draw,angle radius=5pt]{right angle= B--A--C};
	\path pic["$x$", angle eccentricity=2,double,draw,angle radius=9pt]{angle= B--C--A};
	\path pic["$y$", angle eccentricity=2,draw,angle radius=9pt]{angle= A--B--C};
	\end{tikzpicture}
	}
	\loigiai{
	Do góc $x$ và góc $y$ là hai góc phụ nhau nên $\sin y=\cos x\approx 0{,}78$ và $\tan y=\cot x\approx 1{,}25$.
	}
\end{bt}
%%==========Bài 13
\begin{bt}
	\immini{Hình bên mô tả một chiếc thang có chiều dài $AB=4$ m được đặt dựa vào tường, khoảng cách từ chân thang đến chân tường là $BH=1{,}5$ m. Tính góc tạo bởi cạnh $AB$ và phương nằm ngang trên mặt đất (làm tròn kết quả đến hàng đơn vị của độ).}{
	\begin{tikzpicture}[xscale=0.8, yscale=0.6, font=\footnotesize, >=stealth,scale=0.7]
	\path
	(-1,0) coordinate (P)+(-80:1) node{ }
	(P)+(60:5) coordinate (Q)
	(P)+(-30:0.2) coordinate (P1)+(-90:.3) node{$B$}
	(Q)+(-30:0.2) coordinate (Q1)+(10:.37) node{$A$}
	(P)+(150:0.9) coordinate (PP)
	(Q)+(150:0.9) coordinate (QQ)
	(P1)+(150:0.9) coordinate (PP1)
	(Q1)+(150:0.9) coordinate (QQ1)
	(4,0) coordinate (SD)
	(SD)+(150:10) coordinate (ST)
	(Q1)+(-90:1) coordinate (QT)
	(Q)+(0:1) coordinate (Q2)
	(intersection of SD--ST and Q1--QT) coordinate (C)+(-90:.3) node{$H$}
	;
	\draw (-2,-0.5) rectangle (2,5.5);
	\clip (-2,-0.5) rectangle (2,5.5);
	\draw (SD)--(ST) (P1)--(C)--(Q1);
	\draw (Q1)--++(-30:2) (Q1)--++(150:3);
	\draw (Q2)--++(-30:2) (Q2)--++(150:3);
	\pic[draw,angle radius=10]{right angle=P1--C--Q1};
	\foreach \b in {0.1,0.2,...,0.9} {
	\path 
	($(P)!\b!(Q)$) coordinate (m)
	(m)+(50:0.2)coordinate (m1)
	(m)+(150:0.7)coordinate (m2)
	(m1)+(150:0.7)coordinate (m3);
	\draw[fill=cyan!50](m)--(m1)--(m3)--(m2)--(m);}
	\draw[fill=cyan] (P)--(Q)--(Q1)--(P1)--(P);
	\draw[fill=cyan] (PP)--(QQ)--(QQ1)--(PP1)--(PP);
	\end{tikzpicture}}
	\loigiai{
	Ta có, góc tạo bởi cạnh $AB$ và phương nằm ngang trên mặt đất là $\widehat{A B H}$.\\ Xét tam giác $A B H$ vuông tại $H$, ta có $\cos \widehat{A B H}=\dfrac{B H}{A B}=\dfrac{1{,}5}{4}=0{,}375$.\\
	Vậy $\widehat{A B H} \approx 68^{\circ}$.
	}
\end{bt}
%%==========Bài 14
\begin{bt}
	\immini{Treo quả cầu kim loại nhỏ vào giá thí nghiệm bằng sợi dây mảnh nhẹ không dãn. Khi quả cầu đứng yên tại vị trí cân bằng, dây treo có phương thẳng đứng. Kéo quả cầu khỏi vị trí cân bằng một đoạn nhỏ rồi buông ra thì quả cầu sẽ chuyển động qua lại quanh vị trí cân bằng. Khi kéo quả cầu khỏi vị trí cân bằng, giả sử tâm $A$ của quả cầu cách $B$ một khoảng $AB=60$ cm và cách vị trí cân bằng một khoảng $AH=20$ cm (Hình 9). Tính số đo góc $\alpha$ tạo bởi sợi dây $BA$ và vị trí cân bằng (làm tròn kết quả đến hàng đơn vị của độ).}{\begin{tikzpicture}[xscale=0.75,transform shape,scale=.7]
	\definecolor{burlywood}{rgb}{0.87, 0.72, 0.53}%da tay
	\definecolor{blanchedalmond}{rgb}{1.0, 0.92, 0.8}%màu móng
	\definecolor{apricot}{rgb}{0.98, 0.81, 0.69}%màu gỗ bút
	\definecolor{ao(english)}{rgb}{0.0, 0.5, 0.0}%màu bút
	\definecolor{dimgray}{rgb}{0.41, 0.41, 0.41}%màu đinh
	\clip (-2.5,.5) rectangle (3.5,6.5);
	\tikzset{tay_ve/.pic={
	\def\T0{ 
	(-1.75,.65)
	..controls +(-85:.7) and +(150:.2) ..(-1.15,-.3)--(-1,-.13)
	..controls +(-60:.35) and +(-40:.25) ..(-.72,0.05)
	..controls +(-35:.55) and +(-40:.35) ..(-.5,0.3)
	..controls +(80:2) and +(60:1) ..cycle
	;
	}
	\fill[burlywood!85] \T0;
	\draw \T0;
	\draw (-1.3,.6)
	..controls +(45:.1) and +(170:.2) ..(-.7,.85)
	(-1.3,.4)
	..controls +(45:.1) and +(170:.1) ..(-.7,.65);
	\def\T1{ 
	(-.8,2)
	..controls +(-130:.65) and +(125:1.2) ..(-1.6,.32)
	..controls +(-40:.2) and +(125:.2) ..(-1.2,-.28)
	..controls +(-20:.5) and +(-35:.2) ..(-1.3,.55)
	..controls +(80:.15) and +(-75:.15) ..(-1.32,.9)%1
	..controls +(-10:.05) and +(-155:.05) ..(-.9,1.2)%--cycle
	;
	}
	\draw (-1.32,.9)%1
	..controls +(150:.05) and +(-35:.05) ..(-1.45,1);
	\fill[burlywood] \T1;
	\draw \T1;
	\def\T2{ 
	(-.9,1.2)
	..controls +(-40:.1) and +(-150:.2) ..(-.45,1)
	..controls +(-100:.4) and +(80:.45) ..(-1.25,0)
	..controls +(-80:.6) and +(-170:.45) ..(.1,.6)
	..controls +(10:.3) and +(-130:.5) ..(.8,1)
	..controls +(50:.3) and +(-150:.3) ..(1.3,1.5)--(.7,2.2)
	..controls +(-150:.1) and +(35:.1) ..(.4,2)
	..controls +(-150:.1) and +(-10:.1) ..(-.8,2)
	;
	}
	\fill[burlywood] \T2;
	\draw \T2;
	%móng tay
	\draw[fill=blanchedalmond] (-1.05,.27)
	..controls +(10:.05) and +(85:.05) .. (-.85,.05)
	..controls +(-120:.45) and +(-145:.45) .. cycle;
	\draw[ball color=red,opacity=1,draw=black] (-1.17,.2)
	..controls +(-120:.25) and +(-152:.4) .. (-.9,-.1)
	..controls +(-100:.15) and +(-20:.3) .. (-1.3,-.37)
	..controls +(160:.3) and +(-180:.45) .. cycle
	;
	}}
	\fill[gray!20!black,draw=black] (-2.5,1.1)--(3,1.1)--(3,.7)--(-2.5,.7)--cycle;
	\fill[burlywood,draw=black] (-2,2.1)--(3.3,2.1)--(3,1.1)--(-2.5,1.1)--cycle;
	\fill[gray!70!black,draw=black] (3.3,1.7)--(3.3,2.1)--(3,1.1)--(3,.7)--cycle;
	\draw[dashed] (.55,2.2) circle (1.4mm);
	\path (.55,6) coordinate (B)
	(-.8,2.5) coordinate (A)
	(.55,2.5) coordinate (H)
	;
	\fill[gray!90!black,draw=black] (.85,1.4)--(.9,1.8)--(2.1,1.8)--(2.2,1.4)--cycle;
	\fill[gray!50!black,draw=black] (.85,1.2)--(.85,1.4)--(2.2,1.4)--(2.2,1.2)--cycle;
	\fill[gray!50!black,draw=black] (-2.5,1.1)--(3,1.1)--(3,.7)--(-2.5,.7)--cycle;
	\draw[gray!70!black,line width=.5mm] (.2,6)--(2,6);
	\draw[gray!70!black,line width=.5mm] (1.5,6.8)--(1.5,1.65);
	\draw[dashed] (A)--(H)--(B);
	\draw[blue,line width=.3mm] (A)--(B);
	\node at (A) [below]{\tiny $A$};
	\node at (B) [below left]{\tiny $B$};
	\node at (H) [right]{\tiny $H$};
	\node at (.8,5) [left,rotate=90]{\tiny Vị trí cân bằng};
	\node at (H) [right]{\tiny $H$};
	\node at (-1.1,4) [right]{\tiny $60$ cm};
	\node at (-.5,2.3) [right]{\tiny $20$ cm};
	\draw pic["\tiny $\alpha$", draw=black, angle eccentricity=1.5, angle radius=.2cm, color=blue]
	{angle=H--A--B}; 
	\draw pic[draw=black, angle eccentricity=2, angle radius=0.15cm]
	{right angle=B--H--A};	
	\path %(0,0)pic[scale=1]{tay_ve};
	(-1.25,2.25)pic[scale=.42,yscale=-1,rotate=140]{tay_ve};
	\end{tikzpicture}}
	\loigiai{
	Xét tam giác $ABH$ vuông tại $H$, ta có:
	$
	\sin \alpha=\dfrac{A H}{A B}=\dfrac{20}{60}=\dfrac{1}{3}.
	$\\
	Do đó $\alpha \approx 19^{\circ}$.\\
	Vậy góc $\alpha$ tạo bởi sợi dây $BA$ và vị trí cân bằng có số đo khoảng $19^{\circ}$.
	}
\end{bt}
%%==========Bài 15
\begin{bt}
	Cho tam giác $ABC$ vuông tại $A$ có $\widehat{C} = 45^\circ$ và $AB=c$. Tính $BC$ và $AC$ theo $c$.
	\loigiai{
	\immini{
	Xét $\Delta ABC$ vuông tại $A$ có $\sin C = \dfrac{AB}{BC}$\\
	Suy ra $BC= \dfrac{AB}{\sin C} = \dfrac{c}{\sin 45^\circ} = c\sqrt{2}$.\\
	Xét $\Delta ABC$ vuông tại $A$ có $\widehat{B} + \widehat{C} = 90^\circ$.\\
	Do đó $\widehat{B} = \widehat{C} = 45^\circ$.\\
	Suy ra $\Delta ABC$ vuông cân tại $A$.\\
	Suy ra $AC=AB=c$.
	}{
	\begin{tikzpicture}[>=stealth,line join=round,line cap=round,font=\footnotesize,scale=0.5]
	\path 
	(0,0) coordinate[label=left:$A$] (A)
	(5,0) coordinate[label=right:$B$] (B)
	(0,5) coordinate[label=left:$C$] (C);
	\draw (A)--(B)--(C)--cycle;
	\pic[draw, angle radius=3mm]{right angle=B--A--C};
	\pic[draw, angle radius=3mm]{angle=A--C--B} node[below] at ($(C) - (125:0.9)$){$45^\circ$};
	\end{tikzpicture}}}
\end{bt}
%%==========Bài 16
\begin{bt}
	\immini{Tìm chiều cao của tháp canh trong hình bên (kết quả là tròn đến hàng phần trăm).}{
	\begin{tikzpicture}[declare function={a=2.3;},>=stealth, line join=round,line cap=round,font=\footnotesize,scale=0.65]
	\path (0:0) coordinate (B)
	(180:a) coordinate (C)
	(C)+(60:2*a) coordinate (A)
	($(B)!0.5!(C)$) coordinate (D)
	% khai báo vẽ tháp canh
	(180:0.2*a) coordinate (T) %node {T}
	(0:0.2*a) coordinate (D) %node {D}
	(90:1.1*a) coordinate (X) %node {x} % nháp
	(X)+(180:0.13*a) coordinate (E) %node {E}
	(X)+(0:0.13*a) coordinate (F) %node {F}
	% thân chòi
	(X)+(180:0.2*a) coordinate (E1) %node {E1}
	(X)+(0:0.2*a) coordinate (F1) %node {F1}
	(X)+(90:0.45*a) coordinate (Y) %node {Y}
	(Y)+(180:0.3*a) coordinate (Y1) %node {Y1}
	(Y)+(0:0.3*a) coordinate (Y2) %node {Y2}
	($(Y1)!(E1)!(Y2)$) coordinate (E2) %node {E2}
	($(Y1)!(F1)!(Y2)$) coordinate (F2) %node {F2}
	%=========================
	% lang can
	(X)+(180:0.26*a) coordinate (H1) %node {h1}
	(X)+(0:0.26*a) coordinate (H2)% node {h2}
	(H1)+(90:0.1*a) coordinate (H3)% node {h3}
	($(H2)+(H3)-(H1)$) coordinate (H4)% node {h4}
	;
	% thanh chéo ,0.6/T6/X6,0.7/T7/X7,0.8/T8/X8,0.9/T9/X9
	\foreach \i/\t/\x in {0.2/T1/X1,0.4/T2/X2,0.6/T3/X3,0.8/T4/X4}{
	\path ($(E)!\i!(T)$) coordinate (\t);
	\path ($(F)!\i!(D)$) coordinate (\x);
	\draw [thick,black!60!red] (\t)--(\x);
	}
	\draw[thick,black!60!red] (T1)--(X2) (X1)--(T2) (T3)--(X4) (X3)--(T4);
	%==================================
	% chân chòi
	\filldraw[red,opacity=0.65] (A)--(Y1)--(Y2)--cycle;
	\filldraw[orange!70!green,opacity=0.65] (E2)--(E1)--(F1)--(F2)--cycle;
	%==========================
	% vẽ lang can
	\foreach \i/\t/\x in {0/A0/B0,0.2/A1/B1,0.4/A2/B2,0.6/A3/B3,0.8/A4/B4,1/A5/B5}{
	\path ($(H1)!\i!(H2)$) coordinate (\t);
	\path ($(H3)!\i!(H4)$) coordinate (\x);
	\draw[thick,brown] (\t)--(\x);
	}
	\draw[thick,brown] (H1)--(H2)--(H4)--(H3)--cycle;
	%=============
	\draw [thick,black!60!red] 
	(T)--(E)
	(D)--(F)
	;
	%cửa sổ
	\path 
	($(X)!0.5!(F2)$)+(90:0.04*a) coordinate (M) %node {m}
	(M)+(90:0.06*a) coordinate (M1)
	(M)+(0:0.06*a) coordinate (M2)
	(M)+(-90:0.06*a) coordinate (M3)
	(M)+(180:0.06*a) coordinate (M4)
	;
	\draw [blue!70] (M2) rectangle (M1)
	(M3) rectangle (M2)
	(M3) rectangle (M4)
	(M4) rectangle (M1)
	;
	%====================
	\path pic[draw,angle radius=5pt]{right angle= A--B--C};
	\path pic["$60^\circ$", angle eccentricity=2.2, draw,angle radius=8pt]{angle= B--C--A};
	\draw[thick,blue] (A)--(B) 
	--(C) node[midway,below] {$5{,}8$ m}
	--cycle 
	;
	\foreach \t/\g in {A/90,B/-90,C/-90}{
	\path (\t) node[shift={(\g:7pt)}]{$ \t$};
	}
	\end{tikzpicture}}
	\loigiai{
	Xét $\triangle ABC$ vuông tại $B$, ta có\\	
	$\tan C=\dfrac{AB}{CB}$, suy ra $AB=CB\tan C$ hay $AB=5{,}8\cdot\tan 60^\circ=5{,}8\cdot\sqrt{3}\approx 10{,}05 $ (m).\\
	Vậy chiều cao của tháp canh gần bằng $10{,}05$ mét.
	}
\end{bt}
%==========Bài 17
\begin{bt}%[9H1B2]
	Dựng góc $\alpha$, biết $\sin \alpha = 0,25$.
	\loigiai{
	\immini{
	Ta có $0,25 = \dfrac{1}{4}$.
	\begin{itemize}
	\item Dựng góc vuông $xOy$;
	\item Trên cạnh $Ox$ đặt $OA = 1$;
	\item Dựng đường tròn $(A;\ 4)$ cắt cạnh $Oy$ tại $B$.
	\end{itemize}
	Khi đó $\widehat{ABO} = \alpha \ \left( \mbox{vì }\sin \alpha = \dfrac{OA}{AB} = \dfrac{1}{4} \right)$.
	}{
	\begin{tikzpicture}[line join = round, line cap = round,>=stealth,font=\footnotesize,scale=1.3] 
	%
	\tkzDefPoints{0/0/O, 0/1/A, 1/0/E, 0/-3/S} 
	\tkzInterLC(O,E)(A,S) \tkzGetPoints{D}{B}
	\tkzDrawArc[R,delta=20,dashed,color=black](O,1)(80,90)
	\pgfresetboundingbox
	\tkzDrawArc[R,dashed,color=black](A,4)(340,355)
	\draw[->] (0,0)--(5,0) node[below] at (5,-0.2){$x$};
	\draw[->] (0,0)--(0,2) node[right] at (0.2,2){$y$};
	\tkzDrawSegments(A,B)
	\tkzMarkRightAngles[size=0.2](A,O,B)
	\tkzLabelSegments[color=black,left](O,A){$1$}
	\tkzLabelSegments[color=black,above](A,B){$4$}
	\tkzMarkAngles[size=.9,arc=l](A,B,O)
	\tkzLabelAngles[pos=1.2](A,B,O){$\alpha$}
	\tkzDrawPoints[fill=black](A,B) 
	\tkzLabelPoints[below left](O) 
	\tkzLabelPoints[above left](A)
	\tkzLabelPoints[below right](B) 
	\end{tikzpicture}
	}
	}
\end{bt}
%%==========Bài 18
\begin{bt}%[9H1B2]
	Dựng góc $\alpha$, biết $\cos \alpha = 0,75$.
	\loigiai
	{
	Ta có $0,75 = \dfrac{3}{4}$.
	\immini
	{
	\begin{itemize}
	\item Dựng góc vuông $xOy$;
	\item Trên cạnh $Oy$ đặt $OB = 3$;
	\item Dựng đường tròn $(B;\ 4)$ cắt cạnh $Ox$ tại $A$.
	\end{itemize}
	Khi đó $\widehat{ABO} = \alpha \ \left( \mbox{vì }\cos \alpha = \dfrac{OB}{AB} = \dfrac{3}{4} \right)$.
	}
	{
	\begin{tikzpicture}[line join = round, line cap = round,>=stealth,font=\footnotesize,scale=0.9] 
	%
	\tkzDefPoints{0/0/O, 3/0/B, 0/1/E, -1/0/S} 
	\tkzInterLC(O,E)(B,S) \tkzGetPoints{A}{D}
	\tkzDrawArc[R,dashed,color=black](O,3)(-10,10)
	\pgfresetboundingbox
	\tkzDrawArc[R,dashed,color=black](B,4)(130,145)
	\draw[->] (0,0)--(4,0) node[below] at (4,-0.2){$x$};
	\draw[->] (0,0)--(0,3) node[above] at (0.2,3){$y$};
	\tkzDrawSegments(A,B)
	\tkzMarkRightAngles[size=0.2](A,O,B)
	\tkzLabelSegments[color=black,below](O,B){$3$}
	\tkzLabelSegments[color=black,above](A,B){$4$}
	\tkzMarkAngles[size=.5,arc=l](A,B,O)
	\tkzLabelAngles[pos=0.7](A,B,O){$\alpha$}
	\tkzDrawPoints[fill=black](A,B) 
	\tkzLabelPoints[below left](O) 
	\tkzLabelPoints[above left](A)
	\tkzLabelPoints[below right](B) 
	\end{tikzpicture}
	}	
	}
\end{bt}
%%==========Bài 19
\begin{bt}
	Tính giá trị của mỗi biểu thức sau:
	\begin{listEX}[3]
	\item $\sin^2 45^{\circ}+\cos ^2 45^{\circ}$;
	\item $\tan 30^{\circ} \cdot \cot 30^{\circ}$;
	\item $\dfrac{\sin 30^\circ\cdot\cos 60^\circ}{\tan 45^\circ}$
	\end{listEX}
	\loigiai{
	\begin{listEX}
	\item $\sin^2 45^{\circ}+\cos^2 45^{\circ}=\left(\dfrac{\sqrt{2}}{2}\right)^2+\left(\dfrac{\sqrt{2}}{2}\right)^2=\dfrac{1}{2}+\dfrac{1}{2}=1$.
	\item $\tan 30^{\circ} \cdot \cot 30^{\circ}=\dfrac{\sqrt{3}}{3} \cdot \sqrt{3}=1$.
	\item $\dfrac{\sin 30^\circ\cdot\cos 60^\circ}{\tan 45^\circ}=\dfrac{\dfrac{1}{2}\cdot\dfrac{1}{2}}{1}=\dfrac{1}{4}$.
	\end{listEX}
	}
\end{bt}	
%%==========Bài 20
\begin{bt}
	Tính giá trị của các biểu thức sau
	\begin{enumEX}{2}
	\item $A=\dfrac{2\cos 45^\circ}{\sqrt{2}}+\sqrt{3}\tan 30^\circ$;
	\item $B=\dfrac{2\sin 60^\circ}{\sqrt{3}}-\cot 45^\circ$.
	\end{enumEX}
	\loigiai{
	\begin{listEX}
	\item Ta có $A=\dfrac{2\cos 45^\circ}{\sqrt{2}}+\sqrt{3}\tan 30^\circ=\dfrac{2\cdot\dfrac{\sqrt{2}}{2}}{\sqrt{2}}+\sqrt{3}\cdot \dfrac{\sqrt{3}}{3}=1+1=2$;
	\item Ta có $B=\dfrac{2\sin 60^\circ}{\sqrt{3}}-\cot 45^\circ=\dfrac{2\cdot\dfrac{\sqrt{3}}{2}}{\sqrt{3}}-1=1-1=0$.
	\end{listEX}
	}
\end{bt}
%%==========Bài 21
\begin{bt}%[9H1Y2]
	Tính giá trị của biểu thức
	\begin{enumEX}{2}
	\item $M= 4\cos^245^{\circ} + \sqrt{3}\cot30^{\circ} - 16\cos^360^{\circ}$;
	\item $N =\dfrac{2\sin 30^{\circ} - \sin 60^{\circ}}{\cos^230^{\circ} - \cos 60^{\circ}}$.
	\end{enumEX}
	\loigiai
	{
	\begin{enumEX}{2}
	\item \allowdisplaybreaks $\begin{aligned}[t]
	M &= 4\cos^245^{\circ} + \sqrt{3}\cot30^{\circ} - 16\cos^360^{\circ}\\
	&= 4\cdot\left(\dfrac{\sqrt{2}}{2}\right)^2 + \sqrt{3}\cdot\sqrt{3} - 16\cdot\left(\dfrac{1}{2}\right)^3\\
	&= 2 + 3 - 2 = 3.
	\end{aligned}$
	\item \allowdisplaybreaks $\begin{aligned}[t]
	N &=\dfrac{2\sin 30^{\circ} - \sin 60^{\circ}}{\cos^230^{\circ} - \cos 60^{\circ}}\\
	&= \dfrac{2\cdot\dfrac{1}{2} - \dfrac{\sqrt{3}}{2}}{\left(\dfrac{\sqrt{3}}{2}\right)^2 - \dfrac{1}{2}}=\dfrac{1 - \dfrac{\sqrt{3}}{2}}{\dfrac{3}{4} - \dfrac{1}{2}}= 4 - 2\sqrt{3}.
	\end{aligned}$
	\end{enumEX}
	}
\end{bt}
%%==========Bài 22
\begin{bt}
	So sánh
	\begin{enumEX}{4}
	\item $\sin 25^\circ$ và $\cos 65^\circ$;
	\item $ \cos 25^\circ$ và $\sin 65^\circ$;
	\item $\tan 25^\circ$ và $\cot 65^\circ$;
	\item $\cot 25^\circ$ và $\tan 65^\circ$.
	\end{enumEX}
	\loigiai{
	Ta có
	\begin{enumEX}{2}
	\item $\sin 25^\circ=\cos (90^\circ-25^\circ)=\cos 65^\circ$;
	\item $\cos 25^\circ=\sin (90^\circ-25^\circ)=\sin 65^\circ$;
	\item $\tan 25^\circ=\cot (90^\circ-25^\circ)=\cot 65^\circ$;
	\item $\cot 25^\circ=\tan (90^\circ-25^\circ)=\tan 65^\circ$.
	\end{enumEX}
	}
\end{bt}
%%==========Bài 23
\begin{bt}
	\begin{listEX}
	\item So sánh $\sin 72^\circ$ và $\cos 18^\circ$; $ \cos 72^\circ$ và $\sin 18^\circ$; $\tan 72^\circ$ và $\cot 18^\circ$.
	\item Cho biết $\sin 18^\circ\approx 0{,}31$; $\tan 18^\circ\approx 0{,}32$. Tính $\cos 72^\circ$ và $\cot 72^\circ$.
	\end{listEX}
	\loigiai{
	\begin{listEX}
	\item 
	$\sin 72^\circ=\cos (90^\circ-72^\circ)=\cos 18^\circ$;\\
	$\cos 72^\circ=\sin (90^\circ-72^\circ)=\sin 18^\circ$;\\
	$\tan 72^\circ=\cot (90^\circ-72^\circ)=\cot 18^\circ$.
	\item
	$\sin 18^\circ=\cos 72^\circ\approx 0{,}31$;\\
	$\tan 18^\circ=\cot 72^\circ\approx 0{,}32$.
	\end{listEX}
	}
\end{bt}
%%==========Bài 24
\begin{bt}%[9H1B2]
	Sắp xếp các tỉ số lượng giác sau theo thứ tự tăng dần
	\begin{enumEX}{2}
	\item $\cot40^{\circ},\sin 40^{\circ},\cot43^{\circ},\tan42^{\circ}$;
	\item $\tan52^{\circ},\cot63^{\circ},\tan72^{\circ},\cot31^{\circ},\sin 27^{\circ}$.
	\end{enumEX}
	\loigiai
	{
	\begin{listEX}
	\item Ta có $\cot40^{\circ}=\tan50^{\circ}$; $\cot43^{\circ}=\tan47^{\circ}$.\\
	Vì $\sin 40^{\circ}<\tan40^{\circ}<\tan42^{\circ}<\tan47^{\circ}<\tan50^{\circ}$ nên $\sin 40^{\circ}<\tan42^{\circ}<\cot43^{\circ}<\cot40^{\circ}$.
	\item Ta có $\cot63^{\circ}=\tan27^{\circ}$; $\cot31^{\circ}=\tan59^{\circ}$.\\
	Vì $\sin 27^{\circ}<\tan27^{\circ}<\tan52^{\circ}<\tan59^{\circ}<\tan72^{\circ}$ nên $\sin 27^{\circ}<\cot63^{\circ}<\tan52^{\circ}<\cot31^{\circ}<\tan72^{\circ}$.
	\end{listEX}
	}
\end{bt}